\section{Problem Formulation}
\label{sec:formulation}

We study when to compile GUI procedures that agents execute repeatedly while limiting cumulative cost relative to continued agent service.
Tasks in one family share a GUI procedure but differ in their inputs, such as event titles and dates.

We distinguish two ways to serve a task.
In \emph{agent execution}, the model selects GUI actions step by step, with calibrated mean cost $c$.
In \emph{program execution}, a model call extracts the task inputs, or \emph{binding}, and a parameterized program then executes the shared procedure without further model calls.
Its model cost $c^{\mathrm{prog}}$ is the extraction cost.
We use weighted tokens: one unit costs as much as one uncached input token for that model, with cached-input and output tokens weighted by their relative prices.
Appendix~\ref{app:accounting} gives the accounting formula and calibrations.

A compilation attempt uses $k$ completed agent traces.
Let $p$ be its verification probability, $C$ the conditional mean successful compilation cost, and $C^{\mathrm{fail}}$ the conditional mean failed compilation cost.
A failed attempt pays its cost without producing a program.
In the online model, these calibrated costs become fixed charges for the corresponding outcomes.
Historical means alone do not bound future real bills.

Without GUI drift, a program run fails with probability $q$ and requires one agent fallback of cost $c$.
GUI drift breaks a program with probability $h$ per use, requiring fallback and invalidation.
Ordinary failure occurs with probability $q$ conditional on no break.
The cost model assumes detected failures and reports task success separately.
Routing costs $\tau_0+m\ell_t$ when $\ell_t$ programs are listed, where $\tau_0$ is the empty-manifest charge and $m$ is the per-entry charge.

\paragraph{Objective.}
We compare PACE with serving every arrival through agent execution, without compiling or reusing programs.
For arrival $t$, define this agent execution reference charge as $a_t=\tau_0+c$ using that family's supplied cost, and let
\begin{equation}
A_t=\sum_{i=1}^{t}a_i.
\label{eq:agent-reference}
\end{equation}
This reference includes program-served arrivals.
It is a model-based comparison, not an extra agent execution observed on each task.
Let $K_t^\pi$ include policy $\pi$'s routing, serving, fallback, and all compilation charges through arrival $t$, including compilation after the final service.
For an unknown stream length $n$, we seek
\begin{equation}
\min_{\pi\in\Pi_{\mathrm{online}}}\mathbb{E}[K_n^\pi]
\quad\text{subject to}\quad
K_t^\pi\leq(1+\epsilon)A_t,\quad t=1,\ldots,n.
\label{eq:objective}
\end{equation}
The constraint holds for every history allowed by the cost model, while the expectation uses the arrival and outcome model evaluated in the experiments.
The class $\Pi_{\mathrm{online}}$ contains policies that serve each arrival through an agent or an available program, using only the current task, supplied parameters, and past observations.
We supply family identities, cost profiles, routing charges, and $q,h$.
The policy estimates future task repetitions and verification success online.
Artifact lookup and manifest removal incur no token charge.
Savings can fund compilation across families, but unobserved future arrivals add no budget.
The margin is $\epsilon=0.25$ in the main evaluation.
This objective constrains cost, without imposing a task-success guarantee.
